\section{Methodology}
\label{sec:methodology}

\subsection{Approach}
\label{sec:method:approach}

The question requires two things usually provided separately: a model whose behaviour
arises from the structure of the decision process rather than from an optimality
assumption, because the phenomenon of interest is the mismatch between a decentralised
generation market and a centrally planned network; and a designed experiment over the
inputs, because a ranking cannot be read off a handful of named scenarios.

We therefore combine a system dynamics model with a Latin-hypercube designed
experiment. The simulation runs monthly from January 2025 to January 2050 (300 months,
integration step \SI{0.25}{\month}), with an hourly dispatch represented inside each
month by one representative day of 24 hourly levels. The experiment comprises 129 runs:
one central case, two one-at-a-time runs per factor, and 120 Monte Carlo samples in
which all four factors vary simultaneously. One-at-a-time runs isolate each factor
against a fixed reference, which makes a tornado bar attributable but says nothing
about interaction; the Monte Carlo ensemble moves everything at once, which the
percentile envelopes and rank correlations require but leaves individual effects to be
inferred statistically. Reporting both, and checking they agree, is stronger than
either alone.

\subsection{Model structure}
\label{sec:method:structure}

Figure~\ref{fig:model} shows the modules, flows and governing parameters. The full
equation set is in Supplementary Material~S1 and the complete parameter set, with the
workbook cell of every value, in Supplementary Material~S2.

% Block diagram of the model: modules, the flows between them, and the parameters
% that govern each one. The four sensitivity knobs of the experiment are drawn in
% the accent colour so the reader can see where in the structure they act.
%
% Laid out on an explicit coordinate grid rather than with relative positioning:
% the boxes carry enough text that `below=of` chains drift and overlap. Outputs sit
% in a bottom row rather than a right-hand column, which is what keeps the figure
% inside the text width.
\begin{figure*}[tb]
\centering
\begin{tikzpicture}[
  x=1cm, y=1cm, scale=0.92, transform shape,
  mod/.style={draw=black!65, rounded corners=2pt, fill=blue!4, align=center,
              text width=26mm, minimum height=13mm, inner sep=3pt, font=\scriptsize},
  pipe/.style={mod, fill=orange!10},
  outp/.style={draw=black!65, rounded corners=2pt, fill=black!5, align=center,
               text width=34mm, minimum height=8mm, inner sep=3pt, font=\scriptsize},
  par/.style={align=left, font=\tiny, text=black!55, text width=29mm, inner sep=1pt},
  knob/.style={draw=red!65!black, fill=red!8, rounded corners=1pt, align=center,
               font=\tiny\bfseries, inner sep=2pt, text=red!55!black, text width=21mm},
  fl/.style={-{Latex[length=1.5mm]}, draw=black!55, semithick},
  fb/.style={-{Latex[length=1.5mm]}, draw=black!45, semithick, dashed},
  lbl/.style={font=\tiny, text=black!55, inner sep=1pt},
]

\def\rA{7.2}   % top row: needs, decision, generation pipeline
\def\rB{3.2}   % bottom row: demand, dispatch, transmission pipeline
\def\rM{5.2}   % transmission expansion, between the two rows
\def\rO{-0.9}  % outputs

% ============ column 1 ======================================================
\node[mod]  (need) at (0,\rA) {\textbf{Availability and needs}\\[1pt]
  reserve margin and energy balance at horizons $f_0 \dots f_9$};
\node[par]  at (-1.5,\rA-1.3) {$\bar m$ desired margin $=0.25$\\
                               $\sigma$ grid security margin $=0.30$};

\node[mod]  (dem) at (0,\rB) {\textbf{Demand}\\[1pt]
  hourly national demand by horizon, split across zones};
\node[par]  at (-1.5,\rB-1.3) {24 hourly levels $\times$ 420 months\\
                               zonal shares by month};
\node[knob] (kdem) at (0,\rB+1.45) {demand growth $g_D$};

% ============ column 2 ======================================================
\node[mod]  (inv) at (4.3,\rA) {\textbf{Investment decision}\\[1pt]
  NPV per zone--technology; logit allocation};
\node[par]  at (2.8,\rA-1.5) {$\Pi^{*}$ threshold $=1.2$;\;
                              $\theta$ scale $=10^{8}$\\
                              $r$ discount rate $=0.10$\\
                              CAPEX, FOM, ENFICC (NREL)};

\node[mod]  (disp) at (4.3,\rB) {\textbf{Dispatch}\\[1pt]
  hourly merit order, one representative day per month};
\node[par]  at (2.8,\rB-1.45) {24 h $\times$ 12 months\\
                               offer price by technology and zone\\
                               hydro offer from reservoir level};
\node[knob] (kfuel) at (4.3,\rB+1.45) {fuel price $c_f$};

% ============ column 3 ======================================================
\node[pipe] (gpipe) at (8.6,\rA) {\textbf{Generation pipeline}\\[1pt]
  licensing $\to$ construction $\to$ commissioning};
\node[par]  at (7.1,\rA-1.4) {$L^{\text{lic}}_{g}=24$ mo\\
                              $L^{\text{con}}_{g}=24$--60 mo\\
                              $L_{g}$ lifetime $=300$ mo};
\node[knob] (kgen) at (8.6,\rA+1.45) {generation lead time $\Delta_g$};

\node[pipe] (tpipe) at (8.6,\rB) {\textbf{Transmission pipeline}\\[1pt]
  tender $\to$ selection $\to$ licensing $\to$ construction};
\node[par]  at (7.1,\rB-1.5) {lines $97$ mo total\\
                              GETs $73$ mo total\\
                              planning cycle $12$ mo (not scaled)};
\node[knob] (ktr) at (8.6,\rB+1.45) {transmission lead time $\Delta_t$};

% ============ column 4 ======================================================
\node[mod]  (tcat) at (12.9,\rM) {\textbf{Transmission expansion}\\[1pt]
  by demand $\kappa_d$; by generation $\kappa_g$; by interconnection};
\node[par]  at (12.95,\rM-1.6) {$F_{\min}=200$ MW, $F_{\max}=500$ MW\\
                                $\alpha$ allowed GET share $=0.30$\\
                                $c_t=176.61$ USD/km-day};

% ============ outputs (bottom row) =========================================
\node[outp] (out1) at (1.1,\rO)  {Installed capacity by technology [GW]};
\node[outp] (out2) at (6.4,\rO)  {Transmission network [km] and corridor capacity [GW]};
\node[outp] (out3) at (11.7,\rO) {Country average price [USD/MWh]};

% ============ flows ========================================================
\draw[fl] (dem)        -- (need);
\draw[fl] (need)       -- (inv);
\draw[fl] (dem.east)   -- (disp.west);
\draw[fl] (disp)       -- (inv);
\draw[fl] (inv)        -- (gpipe);
\draw[fl] (disp.east)  -- node[lbl, above] {zonal deficits} (tpipe.west);
\draw[fl] (gpipe.east) -- ([yshift=2.5mm]tcat.west);
\draw[fl] (tpipe.east) -- ([yshift=-2.5mm]tcat.west);

\draw[fl] (gpipe.west) -- ++(-0.45,0) |- (out1.east);
\draw[fl] (tcat.south) -- ++(0,-4.0) -| (out2.north);
\draw[fl] (disp.south) -- ++(0,-1.9) -| (out3.north);

% ============ the two feedbacks ============================================
\draw[fb] (tcat.east) -- ++(0.6,0) |- node[lbl, pos=0.75, above]
  {network limits commissioning} (gpipe.east);
\draw[fb] (gpipe.north) -- ++(0,2.5) -| node[lbl, pos=0.22, above]
  {expected capacity at $f_k$} (need.north);

\end{tikzpicture}
\caption{Structure of the model. Solid arrows are material and information flows
within a simulation step; dashed arrows are the two feedbacks that generate the
behaviour of interest --- commissioning is constrained by the network that has
actually been delivered, and investment decisions are taken on expected rather than
current capacity. Grey text gives the governing parameters and their central
values; the four boxes in red are the uncertain drivers varied in the designed
experiment of Section~\ref{sec:method:design}.}
\label{fig:model}
\end{figure*}


\subsubsection{Generation: investment, pipeline, commissioning}
\label{sec:method:gen}

Generation capacity is a stock fed by a pipeline in which committed projects pass
through licensing and construction before commissioning, with total lead time
$(L^{\text{lic}}_{z,g} + L^{\text{con}}_{z,g})(1+\Delta_g)$ --- licensing at 24 months
for every technology, construction from 24 months (PV) to 60 (reservoir hydro).
Commissioning is bounded by the network actually delivered: capacity that has finished
construction cannot enter until there is transmission to evacuate it. This is the first
of the two feedbacks drawn dashed in Figure~\ref{fig:model}, and it is what couples the
generation market to the network plan.

Investment is triggered two ways. Under adequacy, a shortfall of the expected reserve
margin below the desired margin $\bar m$ becomes a capacity requirement. Under
profitability, the model computes a net present value per zone--technology over a
25-year life --- assuming the plant is the only entrant and that current plus committed
capacity persists, with CAPEX, fixed O\&M, fuel and firm-capacity parameters from
\citet{nrel2024} --- admits options above a threshold $\Pi^{*}$, and allocates them
across technologies by a logit share with scale $\theta$. That share represents the
dispersion of investor beliefs rather than a single optimising agent, and is the point
at which the model departs most clearly from a GEP formulation. Decisions use
\emph{expected} rather than current state: expected capacity at 3, 5, 7 and 9 years
includes the pipeline, so investors know the construction times they face. This second
feedback prevents over-investment against a deficit already being remedied.

\subsubsection{Transmission: three categories, three triggers}
\label{sec:method:trans}

Transmission expansion is disaggregated because the categories respond to different
drivers and carry different lead times. Network \emph{for demand} follows zonal peak
growth at an empirical intensity $\kappa_d$ and carries no construction delay, since
the model does not represent intra-zonal networks. Network \emph{for generation}
connects each new plant at intensity $\kappa_g$, triggered by the formal identification
of a project at a planning horizon --- structurally \emph{reactive}, responding to
confirmed rather than anticipated projects, which combined with transmission lead times
exceeding generation lead times is a principal source of the bottleneck studied here.
Network \emph{for interconnection} is triggered when a transfer implied by the dispatch
exceeds delivered corridor capability, met by a grid-enhancing technology if the
shortfall is at most a share $\alpha$ of current capability and by a new line
otherwise, with increments bounded by $F_{\min}$ and $F_{\max}$ to represent lumpiness.

Transmission projects pass a four-stage pipeline --- tender, investor selection,
licensing, construction --- totalling 97 months for lines and 73 for GETs, each stage
scaled by $(1+\Delta_t)$. The 12-month planning cycle is deliberately not scaled: the
model uses it as the period of a modulo operation setting when plans are revised.

\subsubsection{Dispatch, prices and distributed resources}
\label{sec:method:dispatch}

Each month the model dispatches 24 hourly levels by merit order, one plant per
zone--technology pair; reservoir hydro offers between a minimum and a scarcity price as
a function of storage level, giving hydrology a price signal rather than merely an
energy constraint. An unconstrained dispatch sizes the transfers the network would
ideally carry and hence detects reinforcement needs; a constrained one respecting
delivered corridor capability yields congestion, zonal prices and the realised mix. The
reported price is the country average tariff over the final twelve months, averaged
because it oscillates with hydrology. An external path solving the constrained dispatch
with PyPSA \citep{brown2018} exists but is unused here.

Distributed PV adoption is a diffusion process driven by the gap between the levelised
cost of distributed generation and the retail tariff, moderated by a potential market
and a social-influence term \citep{oyarzun2026,antoniolli2022,wu2025}. In every run
reported here, system-level storage and grid-enhancing technologies are \emph{disabled}
at the switch level, as in the shipped baseline, while distributed resources remain
endogenous. This bounds the interpretation: the study quantifies uncertainty around the
current policy configuration, not the effect of deploying storage or GETs.

\subsection{Experimental design}
\label{sec:method:design}

Four drivers are varied (Table~\ref{tab:factors}), selected because each is uncertain,
outside the planner's direct control on the relevant horizon, and argued in the
literature to be a principal risk for a system with Peru's characteristics.

\begin{table}[tb]
  \centering
  \caption{Uncertain factors and ranges. Ranges are ranges of admissibility, not
  estimated distributions.}
  \label{tab:factors}
  \small
  \begin{tabular}{@{}llp{56mm}@{}}
    \toprule
    Factor & Range & Model quantity \\
    \midrule
    Demand growth $g_D$ & $\pm20$\,\% &
      multiplier on national demand, equal to 1 in January 2025 and $1+g_D$ in
      January 2050, so only the growth path changes \\
    Generation lead time $\Delta_g$ & $\pm20$\,\% &
      licensing plus construction time of every generation technology \\
    Transmission lead time $\Delta_t$ & $\pm20$\,\% &
      tender, investor selection, licensing and construction of lines and GETs \\
    Fuel price $c_f$ & $\pm15$\,\% &
      fuel cost of every technology that burns fuel \\
    \bottomrule
  \end{tabular}
\end{table}

Two conventions matter. The demand factor is a \emph{growth-path} factor, not a level
shift, so 2025 demand is untouched and only the rate of growth varies --- the correct
construction for an expansion study. Generation lead times are expressed relative to
the \emph{workbook base case} rather than the planned time: the Peruvian data set
carries an ``other delays'' fraction of $-0.2$, so the factor enters as
$(1+\text{base})(1+\text{change})-1$, giving $-0.36$ and $-0.04$ at the ends. This is
the convention used for transmission, which makes the two delay bars comparable.

Distributions are uniform and independent; there is no information justifying another
shape. The Monte Carlo sample is a Latin hypercube \citep{mckay1992} generated from a
fixed seed and written to disk before any simulation, so the experiment is specified
independently of its execution. The realised design has a maximum absolute correlation
between factors of $0.181$ and a mean of $0.076$; near-orthogonality is what makes the
rank correlations in Section~\ref{sec:results:ranking} attributable to single factors.

\subsubsection{Central, best and worst case}
\label{sec:method:scenarios}

The \textbf{central case} is the system run with every factor at its reference value,
the configuration of the input data set as compiled from Peruvian sources. It is not a
forecast and not the ensemble median: it is the reference against which every other run
is measured, and the tornado bars are departures from it. The \textbf{best} and
\textbf{worst} cases are selected from the 120 Monte Carlo runs on a single declared
metric --- the lowest and highest country average price over the final twelve months
--- which keeps the selection auditable, where a weighted composite would make it
depend on weights with no defensible basis. That ``best'' on price need not be best on
every indicator is a result, not a defect
(Section~\ref{sec:results:extremes}).

\subsection{Implementation and validation}
\label{sec:method:validation}

The model is implemented in Vensim DSS \citep{ventana2024} and executed through its DLL
from a Python driver, each uncertain factor entering as a named constant multiplying
the corresponding structural equation with a default value that reproduces the
unmodified model exactly. Model confidence is established by the standard system
dynamics battery \citep{barlas1996,oliva2004,richardson2024} rather than a single fit
statistic --- structure and parameter verification, dimensional consistency, six
extreme-condition tests, behaviour reproduction and behaviour sensitivity --- reported
in full in Supplementary Material~S3. Re-simulating the central case on the
re-published four-knob model returns \SI{63.85}{\giga\watt} and
\SI{46341.23}{\kilo\meter} at January 2050, matching the reference exactly.
